\chapter{Résultats}
\label{chap:resultats}

Ce chapitre livre le \og combien \fg{} de la question du mémoire, sous la discipline posée
aux chapitres~\ref{ch:identifiabilite} et~\ref{chap:partielle} : les niveaux sont
illustratifs, les écarts et les ordres sont ce qui résiste. Les capitaux sont calculés à
l'échelle du secteur financier, l'entité type de $15\,000$~M\euro{} de provisions servant de
repère forfaitaire et de cible à la descente d'échelle du \S\ref{sec:ancrage}, avec
deux périmètres de sévérité : OpRisk (grandes pertes opérationnelles, queue lourde)
et PRC (sévérité plafonnée à $40$~M\euro). Sauf mention contraire, on lit la médiane
et l'intervalle interquartile, la variance de la sévérité étant infinie ($\hat\xi=0{,}595$ sur
OpRisk, $1{,}03$ sur PRC).

\begin{attention}
\textbf{Comment lire les niveaux de ce chapitre.} Les montants qui suivent sont
\emph{illustratifs} et ne doivent pas être lus comme la mesure du capital d'une entité réelle.
Le chapitre~\ref{chap:partielle} le démontre : en ne faisant varier que la loi de sévérité et la
fréquence, c'est-à-dire des \emph{entrées} du modèle qui dépendent d'un périmètre de collecte
non transposable, le niveau se déplace d'un facteur $41$, tandis que les rapports au socle
varient d'au plus $33\,\%$ et que le groupe de tête du classement des piliers ne change pas. Le
tableau~\ref{tab:scretat} en donne d'ailleurs l'illustration immédiate : le même état conforme
vaut $8\,301$~M\euro{} sous OpRisk et $2\,589$~M\euro{} sous PRC en lecture~A, soit un
facteur~$3{,}2$ pour un seul changement de périmètre de sévérité, facteur qui n'est pas une
constante et grandit sur le socle sans contagion du chapitre~\ref{chap:partielle}. Ce chapitre établit donc des écarts, des
rapports et des ordres ; ses conclusions reposent sur eux, jamais sur un niveau.
\end{attention}

\begin{cle}
\textbf{Réconciliation avec le socle du chapitre~\ref{chap:partielle}.} Ce chapitre parle d'un
SCR \emph{conforme} d'environ $5\,900$~M\euro{}, quand le chapitre~\ref{chap:partielle} parle
d'un \emph{socle} de $5\,275$~M\euro{}. Ce ne sont pas deux mesures de la même chose, mais deux
points d'un seul modèle : le socle est pris à contagion nulle ($W=0$), là où l'état
conforme en conserve une. L'ordre socle $<$ conforme $<$ non conforme est donc cohérent, et le
pont terme à terme comme l'écart résiduel de Monte-Carlo sont établis
\matcomp.
\end{cle}

% SOURCES-SCRIPTS: 65 66 12 16 63

\section{SCR par état de conformité}

Le SCR croît de façon monotone de l'état conforme à l'état non conforme
(figure~\ref{fig:multietats}). Le tableau~\ref{tab:scretat} donne les trois lectures emboîtées :
la lecture~A porte sur la fréquence seule ; la lecture~B ajoute la
propagation de la cascade ; la lecture~C ajoute le canal détection.

\begin{table}[htbp]\centering\small
\renewcommand{\arraystretch}{1.25}
\begin{tabular}{@{}l ccc c ccc@{}}
\toprule
& \multicolumn{3}{c}{\textbf{OpRisk} (M\euro)} & & \multicolumn{3}{c}{\textbf{PRC} (M\euro)} \\
\cmidrule(lr){2-4}\cmidrule(lr){6-8}
État & A & B & C & & A & B & C \\
\midrule
Conforme               & 8301  & 6664  & 5932  & & 2589 & 1903 & 1642 \\
Partiellement conforme & 11059 & 9625  & 9625  & & 3653 & 3118 & 3118 \\
Non conforme           & 15074 & 15074 & 16595 & & 5520 & 5520 & 6577 \\
\bottomrule
\end{tabular}
\caption{SCR par état, trois lectures (A~: fréquence ; B~: $+$~propagation ; C~:
$+$~détection). Le canal détection ne joue qu'aux extrêmes (il
% SOURCES-SCRIPTS: 16
abaisse l'état conforme et relève le non conforme), l'état partiel n'étant pas modulé.}
\label{tab:scretat}
\end{table}

\subsection*{Le chiffre de tête : quatre lectures emboîtées et la lecture publiée}
\label{sec:chiffre-de-tete}

Le tableau qui précède n'est pas le seul protocole du mémoire à chiffrer \og l'écart entre l'état
conforme et l'état non conforme \fg{}. Quatre le font, et ils donnent quatre résultats. Aucun
n'est faux : ils ne relâchent pas le même nombre de canaux entre les deux états, et rien
d'autre ne les sépare. Les publier sans le dire reviendrait à laisser le lecteur trancher sans les éléments.

\begin{center}\footnotesize
\renewcommand{\arraystretch}{1.25}
\begin{tabular}{@{}l c r r r@{}}
\toprule
\textbf{Lecture} & \textbf{canaux} & \textbf{conforme} & \textbf{non conf.} & \textbf{facteur} \\
\midrule
Score de conformité seul, $g$ figé au niveau NC & 1 & 7987 & 16\,847 & $2{,}11$ \\
Fréquence $+$ propagation (lecture~B) & 2 & 6664 & 15\,074 & $2{,}26$ \\
$+$ détection (lecture~C) & 3 & 5932 & 16\,595 & $2{,}80$ \\
\textbf{$+$ accumulation P4 (les quatre canaux)} & \textbf{4} & \textbf{6049} & \textbf{20\,188}
& $\mathbf{3{,}34}$ \\
\bottomrule
\end{tabular}
\end{center}

\vspace{4pt}
\noindent En M\euro, périmètre OpRisk. Les niveaux sont ceux des sorties versionnées ; les
\emph{rapports}, eux, sont recalculés à chaque exécution plutôt que
calculés à la lecture, faute de quoi un niveau amont pourrait bouger sans que le facteur suive.
% SOURCES-SCRIPTS: 25 16 43 67

\begin{cle}
\textbf{Le mémoire publie la lecture à quatre canaux : $6\,049 \to 20\,188$~M\euro, soit un
facteur $3{,}34$ et un écart de $14\,139$~M\euro.} Le motif n'est pas qu'elle donne le chiffre
le plus grand, c'est qu'elle est la seule où l'état conforme est conforme sur tous les
canaux que le modèle possède. Dans la première lecture, l'entité dite conforme propage encore
comme une entité défaillante ; dans les deux suivantes, sa détection ou son accumulation tiers
restent au niveau non conforme. Une entité conforme sur un canal et non conforme sur trois autres
n'est pas un état conforme, c'est une \emph{remédiation partielle}, et la nommer conforme
sous-estime l'écart par construction.

\textbf{Le canal de détection y est inclus parce que son effet est résolu} : il pèse
$2\,928 \pm 343$~M\euro{} en valeur de Shapley et
$4\,562 \pm 819$ en marginal de fermeture, signe résolu sur seize graines.
% SOURCES-SCRIPTS: 68

\textbf{Les trois autres lectures ne sont pas retirées}, délibérément : ce sont des
remédiations \emph{partielles}, ce qui est un objet utile en soi, et leur emboîtement chiffre ce
que coûte de ne fermer qu'une partie des canaux.
\end{cle}

\begin{attention}
\textbf{Le facteur se compare d'une lecture à l'autre, l'écart en euros non} : l'écart passe de $8\,860$ à $8\,410$~M\euro{}
entre les deux premières lignes alors que le facteur monte. La cause n'est pas un effet de modèle
mais un changement de base, l'état conforme passant de $7\,987$ à $6\,664$~M\euro{} d'un
protocole à l'autre. Un écart en euros ne se compare qu'à base égale ; un rapport est sans
dimension. C'est le même motif qui fait publier des élasticités plutôt que des plages au
\S\ref{sec:sensibilite-propagation}.

\textbf{Et la monotonie du facteur s'observe, elle ne se teste pas ici.} Les quatre lignes ne
forment pas une expérience contrôlée, leurs protocoles différant aussi par la graine et par la
base. Elle est \emph{cohérente} avec la composition des canaux mesurée sur les seize
configurations, qui partagent le même moteur et les mêmes graines : les canaux s'y multiplient
presque exactement, le produit des quatre facteurs isolés valant $3{,}47$ pour un facteur publié
de $3{,}34$, si bien que l'interaction de $5\,001$~M\euro{} en euros est l'arithmétique d'un
produit plus qu'un effet propre de la cascade (\S\ref{sec:inverse}).
% SOURCES-SCRIPTS: 100
% SOURCES-SCRIPTS: 68
\end{attention}

\begin{figure}[htbp]\centering
\figover{S33_treillis_canaux_3d.png}
\caption{Les seize configurations de canaux, rangées par nombre de canaux relâchés. La hauteur
est le capital, les deux extrémités étant les deux états publiés, $6\,049$ et $20\,188$~M\euro{} ;
tout sommet intermédiaire est une remédiation partielle. Chaque arête monte, ce qui est la
monotonie du \S\ref{sec:inverse} lue d'un coup d'\oe{}il. La couleur est l'écart à l'additivité,
c'est-à-dire le capital observé moins la somme des quatre effets isolés : elle vaut $+5\,001$ au
sommet complet et $-330$ sur la seule paire négative du dossier, propagation et accumulation. Le
troisième axe sépare les configurations de même rang et ne porte aucune grandeur.}
% SOURCES-SCRIPTS: 68 108
\label{fig:treillis-canaux}
\end{figure}

\noindent\textbf{Où chaque lecture continue de servir.} Les trajectoires du
\S\ref{sec:res-traj} et la priorisation du \S\ref{sec:res-prio} restent conduites en
\emph{lecture~B}. Ce n'est pas une indécision résiduelle mais un choix d'objet : elles portent
sur l'\emph{ordre} de remédiation, dont le chapitre~\ref{ch:robustesse} montre qu'il est
invariant aux leviers testés, de sorte que le canal ajouté déplacerait leurs niveaux sans
déplacer leur conclusion. Les rejouer relèverait par ailleurs de la recalibration, que le gel
interdit.

\begin{figure}[htbp]\centering
\figover{S10_scr_multi_etats.png}
\caption{SCR par état de conformité (trois canaux emboîtés) et bascule des états en crise
systémique.}
% SOURCES-SCRIPTS: 16
\label{fig:multietats}
\end{figure}

Le surcoût de non-conformité $\Delta_{\text{DORA}}=\mathrm{SCR}(\text{NC})-\mathrm{SCR}(\text{C})$,
à graine Monte-Carlo commune (l'écart ne vient que du changement d'état, non du bruit),
s'obtient par bootstrap deux niveaux (multiplicateurs $+$ sévérité), et $100\,\%$ des tirages
sont positifs.

\begin{table}[htbp]\centering\small
\renewcommand{\arraystretch}{1.25}
\begin{tabular}{@{}l cc@{}}
\toprule
& médiane & IC90\,\% \\
\midrule
$\Delta_{\text{DORA}}$ NC vs C (OpRisk) & 7401 & $[2830\,;38793]$ \\
$\Delta_{\text{DORA}}$ PC vs C (OpRisk) & 2556 & $[963\,;12768]$  \\
$\Delta_{\text{DORA}}$ NC vs C (PRC)    & 3668 & $[3201\,;4094]$  \\
$\Delta_{\text{DORA}}$ PC vs C (PRC)    & 1257 & $[1077\,;1466]$  \\
\bottomrule
\end{tabular}
\caption{Surcoût de non-conformité en euros (M\euro), lecture B (fréquence $+$ propagation),
bootstrap deux niveaux. Côté OpRisk la sévérité de queue est ré-échantillonnée sur les
$91$ excès réels, non tirée dans une approximation
% SOURCES-SCRIPTS: 14
normale~; l'intervalle reste large parce que la queue l'est, non par artefact de méthode.}
% SOURCES-SCRIPTS: 16
\label{tab:delta}
\end{table}

\noindent \textbf{Lecture actuarielle.} L'intervalle très large côté OpRisk
(table~\ref{tab:delta}) n'est pas un défaut
du modèle : c'est la signature d'une queue d'indice $0{,}60$, où le niveau du capital est
intrinsèquement imprécis. L'écart typique (interquartile) est bien plus resserré que
l'intervalle à $90\,\%$, et le \emph{signe} du surcoût, lui, est certain. On retrouve la thèse du
mémoire : le SCR est une distribution, pas un point.

\paragraph{L'intervalle PRC est étroit, et il le reste quand on cesse de croire la conversion
exacte.} La sévérité PRC est dérivée du nombre d'enregistrements compromis par la relation
log-log de \citet{jacobs2014}, dont le résidu vaut $0{,}523$ en logarithme naturel : à volume donné,
le coût réel varie d'un facteur $2{,}4$ autour de la droite à $90\,\%$. Réinjecter ce résidu
% SOURCES-SCRIPTS: 21
enregistrement par enregistrement ne rouvre pourtant presque pas l'intervalle, le bruit se
moyennant dans l'ajustement GPD ; il déplace le \emph{niveau} du surcoût de $7\,\%$, sans
l'élargir \matcomp. L'étroitesse de l'intervalle PRC n'est donc pas un artefact
de la conversion déterministe.

\section{Décomposition par pilier}

En basculant un pilier à la fois de conforme à non conforme, le surcoût se décompose
(figure~\ref{fig:decomp}). Le classement obtenu, $P1>P4>P2>P3>P5$, reproduit exactement
le classement ROOT du modèle qualitatif. La convergence n'est pas une validation
indépendante : les amorces de la cascade sont semées selon ROOT (\S\ref{sec:allocation}), et ROOT
se déduit de la même matrice d'expert que $W$ (\S\ref{sec:branch}). Elle atteste la cohérence
du modèle avec son jugement d'origine.

\begin{figure}[htbp]\centering
\figover{S11_decomp_par_pilier.png}
\caption{Décomposition du $\Delta_{\text{DORA}}$ par pilier et validation croisée avec ROOT.
}
% SOURCES-SCRIPTS: 16b
\label{fig:decomp}
\end{figure}

\begin{cle}
\textbf{Une interaction super-additive, établie là où l'estimateur est fiable.} Les surcoûts
isolés ne somment pas au surcoût total : il reste un terme d'\emph{interaction}, signature de la
cascade (un pilier non conforme amplifie aussi les contagions amorcées ailleurs qui le
traversent), qu'un modèle sans propagation (un LDA à copule) ne peut pas produire.
Sa mesure exige toutefois de la prudence : une interaction est une différence de différences,
l'estimateur Monte-Carlo le plus fragile qui soit en queue lourde. Sur la VaR à $99{,}5\,\%$
OpRisk, son \emph{signe n'est pas résolu} : sur six graines à résolution égale
($60\,000$ années chacune), il va de $-137$ à $+1\,395$~M\euro{} et n'est positif que dans
quatre cas sur six, la graine de référence donnant $+1\,395$~M\euro{}. Un estimateur
qui change de signe avec la graine interdit d'en tirer quoi que ce soit. La super-additivité
est en revanche robuste partout où le bruit de queue ne domine pas : en \emph{perte moyenne}
elle est positive sur les six graines ($+532$ à $+659$~M\euro{} OpRisk, $+827$
côté PRC, soit environ un cinquième du surcoût moyen), et jusque dans la VaR du périmètre PRC,
dont le plafond de sévérité tronque la queue ($+905$~M\euro{}). Conséquence pratique : puisque les
marginaux ne somment pas au total, l'attribution du surcoût aux piliers exige une règle
principielle qui somme exactement : c'est l'objet de la valeur de Shapley
(section~\ref{sec:allocation}). À ne pas confondre avec l'interaction entre les quatre
\emph{canaux} de la non-conformité, qui décompose le même écart sur une autre partition, vaut
$+5\,001$~M\euro{} et dont le signe, lui, \emph{est} résolu : elle est reprise terme par terme au
\S\ref{sec:kpi}. Les deux ne s'additionnent pas.
% SOURCES-SCRIPTS: 16b 20b 68
\end{cle}

\section{Allocation du capital : Shapley et Euler}
\label{sec:allocation}

La décomposition qui précède pose deux questions d'allocation distinctes. \emph{Qui porte le
surcoût~?} appelle une attribution par pilier \emph{source} de la non-conformité, qui doit
composer avec l'interaction (les marginaux ne somment pas au total). \emph{Où loge le
capital~?} appelle une allocation du niveau $\mathrm{SCR}(\text{NC})$ par pilier
\emph{touché}, cohérente au sens où les parts somment au total. La valeur de Shapley répond
à la première, la règle d'Euler à la seconde (%
% SOURCES-SCRIPTS: 20
figure~\ref{fig:allocation}).

\subsection{Shapley : attribuer le surcoût, interaction comprise}

On munit les piliers d'une fonction caractéristique
\[
v(S) \;=\; \mathrm{SCR}\bigl(S \text{ en NC},\ \text{reste C}\bigr)-\mathrm{SCR}(\text{tous C}),
\qquad S \subseteq \{P1,\dots,P5\},
\]
calculée par la cascade sur les $2^5$ configurations à graine Monte-Carlo commune. Le
marginal de la section précédente est $v(\{k\})$~; la valeur de Shapley
\[
\varphi_j \;=\; \sum_{S \subseteq \{P1,\dots,P5\}\setminus\{j\}}
\frac{|S|!\,\bigl(5-|S|-1\bigr)!}{5!}\,\bigl[v(S\cup\{j\})-v(S)\bigr]
\]
moyenne la contribution de $j$ sur tous les ordres d'arrivée. Par efficience,
$\sum_j \varphi_j = v(\text{tous}) = \Delta_{\text{DORA}}$ \emph{exactement} : la règle
redistribue l'interaction que les marginaux isolés laissent de côté, quel que soit son signe.

La construction est celle de \citet{Shapley1953}, dont l'efficience est la propriété
axiomatique ; son emploi en allocation de capital remonte à \citet{denault2001}, et son usage
actuariel sous Solvabilité~II est documenté \citep{Decupere2011}. Les cinq piliers autorisent l'énumération exacte des $5!$
ordres d'arrivée ; les algorithmes par échantillonnage des permutations
\citep{CastroGomezTejada2009, SongNelsonStaum2016}, indispensables en grande dimension, ne
sont donc pas nécessaires ici, et aucune approximation n'entre dans le tableau qui suit.

\begin{table}[htbp]\centering\small
\renewcommand{\arraystretch}{1.25}
\begin{tabular}{@{}l rrr c rrr@{}}
\toprule
& \multicolumn{3}{c}{OpRisk} & & \multicolumn{3}{c}{PRC} \\
\cmidrule{2-4}\cmidrule{6-8}
pilier & marginal & $\varphi_j$ & part & & marginal & $\varphi_j$ & part \\
\midrule
P1 gouvernance & 3481 & 3717 & $34{,}0\,\%$ & & 1375 & 1608 & $32{,}4\,\%$ \\
P4 tiers       & 2658 & 2810 & $25{,}7\,\%$ & & 1075 & 1252 & $25{,}3\,\%$ \\
P2 incidents   & 1634 & 1867 & $17{,}1\,\%$ & &  721 &  920 & $18{,}6\,\%$ \\
P3 tests       & 1094 & 1484 & $13{,}6\,\%$ & &  571 &  738 & $14{,}9\,\%$ \\
P5 partage     &  665 & 1048 &  $9{,}6\,\%$ & &  311 &  438 &  $8{,}8\,\%$ \\
\midrule
somme          & 9531 & 10\,927 & $100{,}0\,\%$ & & 4052 & 4957 & $100{,}0\,\%$ \\
\bottomrule
\end{tabular}
\caption{Valeur de Shapley du surcoût $\Delta_{\text{DORA}}$ (M\euro, NC vs C, trois canaux,
protocole des trois canaux). Les marginaux somment à $9\,531$ (OpRisk), la valeur de
% SOURCES-SCRIPTS: 16b
Shapley à $10\,927$, le surcoût total exactement. Cellules arrondies au million et parts au
dixième de point, échelle à laquelle la colonne somme exactement à cent : c'est la propriété
d'efficacité de la valeur de Shapley, et elle se lit dans le tableau plutôt que de s'y
excuser en note.}
% SOURCES-SCRIPTS: 20
\label{tab:shapley}
\end{table}

Le classement de Shapley reproduit le classement ROOT sur les deux périmètres. Une part de
cette convergence est mécanique : les amorces de la cascade sont semées
proportionnellement à ROOT (parts $30/27/18/15/9\,\%$),
% SOURCES-SCRIPTS: 20
et le canal fréquence d'un pilier
module sa propre part d'amorce. Le test non trivial est la \emph{déformation} : la part de
Shapley de P1 ($34\,\%$) excède sa part d'amorce ($30\,\%$), effet propre de la propagation
(P1 est l'émetteur le plus fort de $W$), et l'interaction redistribuée ($+1\,395$~M\euro{}
OpRisk, $+905$ PRC à cette graine) est intégralement un produit de la cascade, absent d'un
modèle additif.

\subsection{Euler : allouer le niveau, et une limite d'estimation assumée}

La règle d'Euler \citep{tasche1999} alloue $\mathrm{SCR}(\text{NC})$ par
$C_j = \mathbb{E}\bigl[L_j \mid L \ge \VaR_{99{,}5\,\%}\bigr]$ (version TVaR, exacte
par construction), où $L_j$ est la perte du pilier \emph{touché} $j$ dans la cascade. Côté
PRC, l'allocation est stable et presque plate : P2 $23\,\%$, P4 $23\,\%$, P1 $20\,\%$,
P3 $18\,\%$, P5 $17\,\%$, et les parts de capital coïncident avec les parts de perte moyenne
parce que le plafond de sévérité tronque la queue
($\TVaR = 6\,795 \approx \VaR = 6\,573$~M\euro). Côté OpRisk, les parts
\emph{moyennes} sont stables (P2 $23\,\%$, P4 $22\,\%$, P1 $20\,\%$, P3 $18\,\%$,
P5 $16\,\%$), mais l'allocation de la queue extrême ne l'est pas, et la cause en est
théorique : avec $\xi \approx 0{,}6 > 0{,}5$, la sévérité est de variance infinie,
l'estimateur Monte-Carlo de la TVaR converge en régime de loi stable, et le classement
conditionnel bascule d'une résolution à l'autre (vérifié à $150$ et $240$~mille années
simulées). On rapporte donc l'allocation de queue du périmètre PRC, régularisée par le cap,
et l'on tient celle d'OpRisk pour indicative.

\paragraph{La mesure de queue sert ici à \emph{mesurer}, jamais à couvrir.} La distinction est
assez souvent négligée pour qu'il faille l'écrire. Le besoin de capital publié dans ce mémoire
est partout un $\VaR_{99{,}5\,\%}(L)$ centré, conformément au niveau de Solvabilité~II ; aucune
grandeur de couverture n'est une TVaR, ni une CTE, ni un Expected Shortfall. La mesure
de queue intervient à deux titres, et tous deux sont des titres de \emph{diagnostic}. Comme
noyau d'allocation d'abord, dans la règle d'Euler ci-dessus, où elle répartit un niveau déjà
fixé sans le modifier. Comme \emph{indicateur de forme de queue} ensuite : le rapport
$\TVaR_{99{,}5\,\%}(L)/\VaR_{99{,}5\,\%}(L) \approx 2{,}1$ se compare à son asymptote théorique
$1/(1-\xi) = 2{,}5$. L'information est portée par ce rapport, non par les deux montants.
% SOURCES-SCRIPTS: 67
Ce rapport dépend de l'état et il faut
donc dire lequel : $2{,}10$ vaut à l'état \emph{tous non conforme}, celui de cette section, tandis
qu'il vaut $2{,}17$ sur la calibration de référence \matcomp. Les deux
% SOURCES-SCRIPTS: 73
valeurs sont justes, elles ne portent pas sur la même loi de perte, et l'écart entre elles est
sans signification en soi.

Deux raisons de ne pas en faire un niveau de couverture, dont la seconde est décisive. La
première est de proportion : sur un risque aussi large, dont les queues sont plurielles selon la
catégorie d'événement (chapitre~\ref{chap:adaptations}), adosser la couverture à une moyenne de queue
déformerait la distribution et immobiliserait un capital sans commune mesure avec le niveau
réglementaire de référence. La seconde est d'existence : sous la calibration PRC,
$\hat\xi=1{,}033>1$ place la sévérité en régime d'espérance infinie, et la mesure
n'existe alors pas au sens mathématique. Une mesure de couverture qui cesse d'être définie
selon la source de sévérité retenue ne peut pas porter un besoin de capital ; le même objet,
lu comme un diagnostic de queue, reste au contraire parfaitement informatif, puisque sa
divergence \emph{est} l'information.

\begin{cle}
\textbf{Sources et réceptacles.} Les deux allocations ne racontent pas la même chose, et
c'est le résultat : Shapley, vue \emph{source}, concentre le surcoût sur P1 ($34\,\%$), là
où naît la non-conformité~; Euler, vue \emph{réceptacle}, répartit le capital presque
uniformément sur les piliers touchés (P2 et P4 en tête légère). La remédiation se décide
côté source~; le capital, lui, se loge là où les cascades atterrissent.
\end{cle}

\begin{figure}[htbp]\centering
\figover{S15_allocation_shapley_euler.png}
\caption{Allocation du capital DORA (OpRisk). (a)~Shapley contre marginaux : l'interaction
redistribuée. (b)~Euler : parts de capital (queue) contre parts de perte moyenne~; sur ce
périmètre non plafonné, l'allocation fine de la queue extrême reste indicative
($\xi \approx 0{,}6$, variance de sévérité infinie).}
% SOURCES-SCRIPTS: 20
\label{fig:allocation}
\end{figure}

\section{Trajectoire du capital}
\label{sec:res-traj}

Le long de la mise en conformité, le capital décroît de l'état initial vers le SCR conforme, et
le surcoût résiduel $\Delta_{\text{DORA}}(t)$ tend vers zéro (tableau~\ref{tab:traj},
figure~\ref{fig:traj}). Partant de l'état marginal actuel (NC~$35\,\%$, PC~$35\,\%$, C~$30\,\%$),
le capital OpRisk passe de $\sim 10\,750$ à $\sim 7\,040$~M\euro{} sur cinq ans, vers le SCR
conforme ($\sim 6\,925$).

\begin{table}[htbp]\centering\small
\renewcommand{\arraystretch}{1.2}
\begin{tabular}{@{}c ccc@{}}
\toprule
$t$ (ans) & SCR($t$) médiane & bande 90\,\% & $\Delta_{\text{DORA}}(t)$ \\
\midrule
0 & 10752 & $[9092\,;12654]$ & 3827 \\
1 & 9576  & $[7205\,;12040]$ & 2652 \\
2 & 8310  & $[6194\,;11042]$ & 1385 \\
3 & 7575  & $[5924\,;10164]$ & 650 \\
5 & 7041  & $[5857\,;9028]$  & 116 \\
\bottomrule
\end{tabular}
\caption{Trajectoire du capital DORA (M\euro, OpRisk).}
% SOURCES-SCRIPTS: 17
\label{tab:traj}
\end{table}

\begin{figure}[htbp]\centering
\figover{S12_trajectoire_scr.png}
\caption{Trajectoire SCR($t$) et évolution des états de conformité.}
% SOURCES-SCRIPTS: 17
\label{fig:traj}
\end{figure}

\noindent \textbf{Lecture actuarielle.} L'intégrale de la trajectoire, à un coût de capital
près, est le \emph{coût de portage} de la non-conformité pendant la remédiation. La bande reste
large et surtout haute en environnement dégradé : la vraie exposition n'est pas le capital
médian, c'est le risque que la remédiation traîne quand la menace monte, précisément le scénario
où le capital importe le plus.

\section{Priorisation de la remédiation}
\label{sec:res-prio}

Sous budget contraint (un pilier remédié à la fois), l'énumération des $120$ ordres donne
l'ordre optimal $P1\to P4\to P2\to P3\to P5$, égal à l'ordre ROOT et distinct de l'ordre
glouton par valeur d'accélération (figure~\ref{fig:prio}). L'écart de capital intégré entre le
pire ordre et l'optimal atteint $21\,\%$.

\begin{figure}[htbp]\centering
\figover{S13_priorisation.png}
\caption{Valeur d'accélération par pilier et ordre de remédiation optimal.}
% SOURCES-SCRIPTS: 18
\label{fig:prio}
\end{figure}

\begin{cle}
\textbf{Le myope n'est pas l'optimal.} Remédier au coup par coup le pilier au plus fort gain
immédiat (l'ordre glouton) coûte plus cher que l'ordre optimal, parce que les interactions de
cascade imposent de raisonner sur l'horizon entier. L'optimum coïncide avec le classement
ROOT, comme la progéniture (\S\ref{sec:branch}) et la décomposition par pilier ; ces trois
convergences attestent la cohérence du modèle avec le jugement dont il procède, non une
validation indépendante, ROOT entrant dans les amorces.
\end{cle}

\paragraph{La question réciproque donne le classement qu'un plan de
remédiation utilise.} Tout ce qui précède fixe un état et lit un capital. L'exercice de
résistance \emph{inverse} fixe le capital et cherche les états qui le produisent. Un dispositif ORSA emploie un modèle sous
cette forme. La monotonie établie au \S\ref{sec:kpi} le
rend court : l'ensemble des configurations atteignant une cible est croissant, donc entièrement
décrit par ses éléments minimaux, et il n'en compte jamais plus de quatre sur seize. Deux nombres
résument alors tout le treillis. Le canal de fréquence seul atteint $10\,377$~M\euro,
soit $1{,}72$ fois l'état conforme, tandis que les trois autres canaux \emph{réunis} n'atteignent
que $11\,413$, soit $1{,}89$ fois. Toute cible sous le premier seuil se produit donc par un canal
unique, et toute cible au-dessus du second exige la fréquence. Aucun autre canal, seul et
dans sa plage déclarée, n'atteint même $8\,000$~M\euro. La lecture directe hiérarchise les canaux
par leur contribution ; la lecture inverse désigne celui dont la maîtrise \emph{interdit} les
scénarios sévères. L'information n'est pas la même. Détail, frontières par cible et réserves
au \S\ref{sec:inverse}.
% SOURCES-SCRIPTS: 92

\section{Robustesse}

Le niveau du capital bouge fortement selon les leviers non calibrés (l'indice de queue $\xi$
domine, facteur $\sim 6$), mais la direction ($\Delta_{\text{DORA}}>0$) et la
\emph{priorité} (P1 en tête) tiennent sur tous les réglages testés (figure~\ref{fig:robust}).

\begin{figure}[htbp]\centering
\figover{S14_robustesse_multietats.png}
\caption{Robustesse : le classement tient, le niveau non.}
% SOURCES-SCRIPTS: 19
\label{fig:robust}
\end{figure}

\begin{cle}
\textbf{La recommandation ne dépend d'aucun paramètre non mesurable.} L'ordre de remédiation
optimal ne dépend que des SCR par configuration, donc il est invariant aux taux de
transition de la chaîne de Markov, pourtant les paramètres les moins calibrables. Ces taux
fixent le calendrier de la trajectoire, jamais la séquence recommandée. Le verdict d'ensemble
rejoint celui de toute la démarche : le classement est robuste là où le niveau ne l'est pas.
\end{cle}

\begin{attention}
\textbf{Le périmètre exact de cette robustesse.} Les tests ci-dessus font
varier les paramètres du modèle \emph{à matrice $W$ fixée}, celle du classeur qualitatif. Ils
établissent donc que P1 arrive en tête \emph{si l'on admet cette direction}. Dès lors que l'on
admet ne pas connaître la direction et que l'on balaie l'ensemble admissible
(chapitre~\ref{chap:partielle}), P1 et P4 deviennent pratiquement à égalité : P1 arrive
premier dans $43{,}6\,\%$ des configurations, P4 dans $37{,}7\,\%$, et l'écart de regret maximal entre
% SOURCES-SCRIPTS: 30
les deux n'est que de $1{,}5\,\%$. Les deux énoncés ne se contredisent pas, ils portent sur des
sources d'incertitude différentes. Ce qui est robuste dans les deux cadres, c'est que P1 et P4
devancent nettement P2, P3 et P5, et que P5 n'arrive \emph{jamais} en tête.

\medskip
\noindent\textbf{Une troisième source d'incertitude, et elle vise directement le corpus.} Les
rapports post-incident sont des enquêtes officielles, mandatées pour établir une cause racine ;
elles sur-attribuent donc structurellement à la gouvernance, ce que le
chapitre~\ref{ch:identifiabilite} déclare sans pouvoir le lever. On peut néanmoins en chiffrer la
conséquence, en dégradant l'émission de P1 par $W_{1k}\to(1-a)W_{1k}$ et en recalculant l'ensemble
admissible à chaque degré. Le classement tient jusqu'à $a=0{,}25$ et bascule vers P4 à
$a=0{,}50$ : il faut donc supposer la moitié des arêtes issues de la gouvernance
artefactuelles pour renverser la priorité. Le niveau, lui, ne perd que $24{,}8\,\%$ même à
$a=1$, soit $0{,}81$ fois la largeur de la bande d'identification. Le niveau résiste, le
classement non, le seuil est connu.
% SOURCES-SCRIPTS: 64
\end{attention}

\paragraph{Les hypothèses restantes ont leur axe de stress.} Trois entrées ne figuraient pas au
tornado : les coefficients de la conversion Jacobs, la charge systémique $\gamma$ et l'ancrage
des probabilités d'état. Stressées à $\pm 1{,}645$ écart-type, elles déplacent des niveaux sans
jamais toucher la direction ni le classement : le surcoût PRC reste contenu entre $3\,125$ et
$3\,977$~M\euro{} autour du central, $\gamma$ ne commande que la sévérité de la bascule en
crise, et l'ancrage des probabilités d'état déplace le capital espéré actuel mais laisse
inchangés le SCR par état et le surcoût NC contre C, qui lui sont conditionnels. Le détail du
protocole de stress est déporté \matcomp.
% SOURCES-SCRIPTS: 22 67

\section{La sensibilité aux probabilités de propagation, prise pour elle-même}
\label{sec:sensibilite-propagation}

% SOURCES-SCRIPTS: 15 76 30 66 68

Le modèle repose sur des probabilités de propagation qui ne sont pas observées. Cette section
leur consacre un examen séparé, en
distinguant deux questions que l'on confond volontiers : la sensibilité au \emph{niveau} de la
propagation, qui est un facteur d'échelle posé, et la sensibilité à sa \emph{direction}, qui
n'est pas identifiable. La première se stresse, la seconde se borne.

\paragraph{Le niveau de propagation est le levier le moins sensible.} Mesurée en élasticités, à
perturbation relative de $10\,\%$ appliquée à un paramètre à la fois, la sensibilité classe les
leviers comme suit.

\begin{center}\footnotesize
\renewcommand{\arraystretch}{1.15}
\begin{tabular}{@{}l l r@{}}
\toprule
\textbf{Levier} & \textbf{Statut} & \textbf{Élasticité} \\
\midrule
Indice de queue $\xi$ & estimé & $3{,}79$ \\
Échelle $\sigma$ & estimé & $0{,}97$ \\
Taux de dépassement $p_u$ & estimé & $0{,}82$ \\
Fréquence $\lambda$ & estimé & $0{,}45$ \\
Sur-dispersion $\varphi$ & posé & $-0{,}44$ \\
\textbf{Gain de propagation} $\boldsymbol{g}$ & \textbf{posé} & $\mathbf{0{,}37}$ \\
Seuil $u$ seul & choix de méthode & $0{,}03$ \\
\bottomrule
\end{tabular}
\end{center}

\vspace{4pt}
\noindent La propagation a l'élasticité la plus faible des leviers substantiels, $0{,}37$ contre
$3{,}79$ pour l'indice de queue. L'ablation le confirme par un autre chemin : retirer la
propagation fait tomber la VaR de $20\,\%$, retirer la queue lourde de $76\,\%$. La thèse ne
dépend pas des trois valeurs posées de $g$, le capital étant croissant en $g$
(\S\ref{sec:invariance-g}). La direction, elle, ne se stresse pas, elle se borne : à ignorance
totale la largeur de l'ensemble admissible vaut $1\,839$~M\euro, soit $34{,}9\,\%$ du socle, et
elle se paye linéairement avec le degré d'ignorance. Le motif qui écarte les indices de Sobol, la
variance de la charge annuelle n'existant pas, le tornado brut et ses deux défauts, la gradation
de l'ignorance, l'ablation et la comparaison avec le préprint de cascade climatique, qui classe
les leviers à l'inverse parce que sa sévérité est bornée, sont déportés
\matcomp.
% SOURCES-SCRIPTS: 81

\begin{cle}
\textbf{Ce que cette section établit.} La sensibilité aux probabilités de propagation est
\emph{bornée par construction} et non estimée : c'est la différence entre un modèle qui pose une
direction et un modèle qui assume ne pas la connaître. Et elle n'est pas le point faible du
dispositif : le niveau de propagation déplace le capital moins que tout autre levier, tandis que
le seuil de la loi de queue le fait passer à lui seul de $4\,862$ à $12\,944$~M\euro. Toute vigilance sur les
valeurs numériques du modèle doit donc se porter d'abord sur l'ajustement de valeurs extrêmes,
et c'est là que le mémoire concentre ses réserves déclarées.
\end{cle}

\section{Comparaison de modèles : la dépendance n'est pas le mécanisme}
\label{sec:benchmark}

La cascade dirigée remplace ici l'agrégation par copule d'un modèle LDA. Ce que ce
remplacement apporte se mesure en confrontant la cascade à des jumeaux copule construits pour lui
ressembler (figure~\ref{fig:benchmark}, périmètre OpRisk, graines
% SOURCES-SCRIPTS: 23
Monte-Carlo communes). Cette section compare les modèles sur l'axe de la \emph{structure de
dépendance}, à marges fixées ; la comparaison sur l'axe de la \emph{famille de loi} (la bande
de modèle) est conduite en partie~\ref{part:limites}, section~\ref{sec:trois-etages}, et les
deux se complètent.

\textbf{La photo.} À marges par pilier \emph{identiques} (celles de la cascade, état
conforme) et corrélation calée, la copule gaussienne retombe sur le SCR de la cascade
($+0{,}1\,\%$) : en queue lourde ($\xi \approx 0{,}6$), la queue annuelle obéit au principe
de la perte unique dominante et le \emph{niveau} de base n'identifie pas la dépendance. La
Student ($\nu = 4$) le surestime de $17\,\%$ en forçant des co-extrêmes que le mécanisme ne
produit pas ($2{,}28$ piliers extrêmes par année de queue contre $1{,}12$). Le choix de
copule est donc à la fois non identifiable et conséquent~; la cascade n'a pas ce degré de
liberté : sa dépendance est produite par le mécanisme, pas choisie.

\textbf{Et ce résultat a été répliqué ailleurs, sans lien avec ce travail.} L'objection naturelle
serait qu'une coïncidence à $+0{,}1\,\%$ entre une copule gaussienne et un mécanisme de cascade
tient à un réglage particulier du protocole. Le préprint de cascade climatique cité au
chapitre~\ref{ch:positionnement} \citep{ccrn2026} conduit la même expérience, à marges appariées,
entre une gaussienne, une Student et sa propre cascade : il obtient des primes centrales
pratiquement indiscernables et ne voit la séparation apparaître qu'en \emph{queue}. Même
protocole, même conclusion, sur un autre péril et par une autre équipe. Ce n'est donc pas un
artefact du réglage retenu ici, mais une propriété du principe de la perte unique dominante,
qui rend le niveau central aveugle à la structure de dépendance.

Le $+0{,}1\,\%$ est cependant un tirage : sur quatre graines, l'écart moyen vaut $+4{,}8\,\%$
% SOURCES-SCRIPTS: 80
et reste sous son propre bruit, et au-delà de $99\,\%$ seule la Student sort du sien. Le détail
est déporté \matcomp.

\textbf{L'intervention.} Basculer un pilier en non-conformité est une intervention, pas un
conditionnement. Le jumeau LDA-copule exprime les canaux fréquence et détection (les marges
dépendent de l'état de leur pilier), mais trois manques sont structurels.
\emph{(i)}~La propagation n'a nulle part où vivre : le surcoût total ressort à
$6\,760$~M\euro{} contre $10\,927$ pour la cascade, une sous-estimation de $38\,\%$.
% SOURCES-SCRIPTS: 67
\emph{(ii)}~Les marges ne dépendant que de l'état de leur propre pilier et la copule étant
figée, la perte moyenne totale est \emph{exactement} additive : interaction moyenne nulle
par construction ($+0$~M\euro{} mesuré), là où la cascade en produit $+560$~M\euro{} :
l'interaction moyenne mesurée plus haut est une signature falsifiable qu'aucun modèle
copule-marges ne peut reproduire. \emph{(iii)}~La priorité de remédiation impliquée
s'inverse sur P2/P3, et la dépendance ne durcit pas quand la conformité se dégrade, alors
que c'est précisément le canal que DORA fait bouger ($g$ de $0{,}45$ à $0{,}90$).

\begin{cle}
\textbf{Le verdict du benchmark.} La copule peut imiter la photo d'un état~; elle ne
peut imiter ni le surcoût (propagation), ni l'interaction (additivité en moyenne), ni le
durcissement de la dépendance avec la non-conformité. La dépendance n'est pas le
mécanisme : c'en est l'ombre portée.
\end{cle}

\begin{figure}[htbp]\centering
\figover{S18_benchmark_copule.png}
\caption{Benchmark copule (OpRisk). (a)~Intervention : le jumeau copule suit les marges et
sous-estime le surcoût, la cascade suit les sources. (b)~Photo : perte unique dominante en
queue~; la Student force des co-extrêmes que le mécanisme n'a pas.}
% SOURCES-SCRIPTS: 23
\label{fig:benchmark}
\end{figure}

\section{Mise en regard réglementaire}
\label{sec:regard-reglementaire}

La question finale est celle de la portée : de l'effet de la conformité que le modèle
mesure, quelle fraction les cadres existants savent-ils exprimer ? On confronte donc, sur
le même écart de capital entre profil conforme et non conforme, trois dispositifs.
% SOURCES-SCRIPTS: 27

\begin{table}[htbp]\centering\small
\renewcommand{\arraystretch}{1.25}
\begin{tabular}{@{}l r r@{}}
\toprule
\textbf{Cadre} & \textbf{$\Delta$ capital C$\to$NC} & \textbf{part de l'effet cascade} \\
\midrule
Formule Standard (Solvabilité~II) & $0$~M\euro{} & $0\,\%$ \\
LDA-copule (dépendance symétrique) & $6\,760$~M\euro{} & $62\,\%$ \\
Cascade dirigée & $10\,927$~M\euro{} & $100\,\%$ (référence) \\
\bottomrule
\end{tabular}
\caption{Fraction de l'effet DORA que chaque cadre sait exprimer.}
% SOURCES-SCRIPTS: 27 67
\label{tab:benchmark-sf}
\end{table}

\textbf{La Formule Standard est structurellement aveugle}, et la table~\ref{tab:benchmark-sf}
chiffre ce que chaque cadre sait en exprimer. La charge de risque opérationnel
de l'équation~\eqref{reg:eq:scrop}, maximum d'une branche assise sur les primes et d'une branche
assise sur les provisions, à coefficients distincts en vie et en non-vie, dans la limite d'un
plafond assis sur le BSCR, ne dépend que du volume : elle est \emph{identique} pour une entité
exemplaire et pour une entité défaillante sur DORA. Sa réponse à la conformité est nulle par
construction. La copule en capte $62\,\%$ (les canaux fréquence et détection) mais rate la
propagation, faute de direction. Seule la cascade exprime l'effet complet. La valeur du
modèle se lit donc en effet DORA rendu visible, non en un niveau de SCR à
interpréter dans l'absolu.

Sur donnée ouverte, le pilotage descend jusqu'aux sous-piliers de P2 et, plus grossièrement, à
la détection de P3, mais pas jusqu'au pilier des tiers, que seul un registre d'externalisation
peut renseigner \matcomp.

\section{Un ordre de grandeur à l'épreuve du réel}
\label{sec:ancrage}

Le niveau absolu n'ayant jamais été présenté comme une mesure, il reste à le confronter au
réel, non pour le valider mais pour vérifier qu'il ne \emph{peut} pas l'être. Deux
confrontations suffisent.

Contre la Formule Standard, sur la même entité notionnelle, le repère forfaitaire vaut environ
$450$~M\euro{} quand le socle de la cascade au niveau du secteur en vaut $5\,275$ ; contre le
marché, un capital de cet ordre pour une seule entité serait invraisemblable. Au niveau du
sinistre unique, en revanche, le quantile de sévérité ($663$~M\euro) et sa moyenne de queue
($1\,752$~M\euro) encadrent les grandes pertes cyber réelles. Le niveau sectoriel tient donc à la
fréquence d'amorce, et la descente d'échelle le vérifie : à matrice $W$ et moteur inchangés, la
fréquence est lue à la taille de l'entité par une binomiale négative à lien logarithmique
(élasticité $0{,}074$) et la sévérité transposée par son élasticité à la taille ($0{,}087$,
multiplicateur $0{,}854$).

\begin{table}[htbp]\centering
\small
\begin{tabular}{lrrrr}
\toprule
Lecture & $\lambda$ (an$^{-1}$) & sévérité & SCR $99{,}5\,\%$ & Perte moyenne \\
\midrule
Secteur financier mondial & $21{,}56$ & $\times 1$ & $8\,123$~M\euro & $983$~M\euro \\
Seuil \og{}$\geq 10$ événements\fg{} & $0{,}210$ & $\times 1$ & $488$~M\euro & $9{,}7$~M\euro \\
Seuil \og{}toutes firmes\fg{} & $0{,}100$ & $\times 1$ & $229$~M\euro & $4{,}6$~M\euro \\
$\lambda$ lu à la taille de la cible & $0{,}092$ & $\times 1$ & $198$~M\euro & $4{,}3$~M\euro \\
\textbf{Les deux canaux composés} & $\mathbf{0{,}092}$ & $\mathbf{\times 0{,}854}$ &
$\mathbf{169}$~M\euro & $\mathbf{3{,}6}$~M\euro \\
\bottomrule
\end{tabular}
\caption{Le même modèle lu à plusieurs échelles. La matrice $W$ et le moteur de cascade sont
identiques partout ; seules la fréquence d'amorce et l'échelle de sévérité changent. Les deux
dernières lignes sont la lecture cohérente, où la fréquence \emph{et} la sévérité sont prises
à la taille de l'entité visée.}
% SOURCES-SCRIPTS: 60
\label{tab:echelle}
\end{table}

Diviser $\lambda$ par $235$ divise la perte moyenne par $270$ mais le capital par $48$ seulement :
le quantile extrême est porté par le plus grand sinistre, non par le nombre. La coïncidence
apparente avec le repère forfaitaire au seuil de dix événements ($488$ contre $450$~M\euro) ne
survit pas à la correction, et le chiffre d'entité est une bande, $[131{,}5\,;183{,}3]$~M\euro{}
sur les $1\,024$ sommets admissibles. Ce $169$~M\euro{} est une borne supérieure d'ordre de
grandeur (\S\ref{sec:entites-reelles}).
% SOURCES-SCRIPTS: 60 49

\begin{cle}
\textbf{À l'échelle d'entité, ignorer la direction coûte autant qu'ignorer la queue.} La
bande d'identification mesure $51{,}7$~M\euro{} de large. En faisant varier le seul indice de
queue sur son intervalle de confiance à $90\,\%$ ($\xi \in [0{,}304\,;0{,}831]$), le SCR
d'entité va de $140{,}9$ à $191{,}6$~M\euro, soit $51{,}5$~M\euro{} de large. Les deux
incertitudes sont du même ordre, et aucune ne domine l'autre. C'est un argument
supplémentaire pour la spécification de reporting : réduire l'ignorance directionnelle vaut,
en euros de capital, autant qu'affiner l'estimation de la queue.
\end{cle}

\paragraph{Ce que le forfait ne voit pas, à l'échelle d'entité.} En balayant le gain de
propagation sur ses trois valeurs d'état ($0{,}45$, $0{,}68$, $0{,}90$) à fréquence et taille
cibles, le SCR d'entité passe de $116$ à $143$ puis $169$~M\euro{}, soit $+45\,\%$ entre le
profil conforme et le profil défaillant, et ce par le seul canal de contagion, les canaux
fréquence et détection n'étant pas mobilisés ici : l'écart lu est donc un plancher. Sur les
mêmes trois profils, le repère forfaitaire vaut $450$~M\euro{} et ne bouge pas d'un euro, pas
plus que la charge réglementaire elle-même, qui ne dépend d'aucun état. C'est la thèse du
mémoire, lue à l'échelle où une entité la vit.

Ces trois nombres sont un scénario d'écartement, non une mesure : la correspondance
état $\to g$ est posée et non calibrée, faute d'entité publiant son état de conformité. Ce qui
est établi, et démontré au \S\ref{sec:invariance-g}, est plus fort et plus modeste à la fois :
le capital étant croissant en $g$, l'écart existe et garde son signe pour \emph{toute}
correspondance respectant l'ordre, y compris la plus étroite, quand le forfait reste plat dans
tous les cas. Seule l'amplitude dépend de l'écartement retenu.

Les deux canaux de la descente, ce que la correction d'échelle déplace, le rejet d'une mise à
l'échelle proportionnelle et les figures sont déportés \matcomp.

\section{Quatre entités réelles, et la borne inférieure de la méthode}
\label{sec:entites-reelles}

Tout ce qui précède porte sur une entité \emph{notionnelle}, posée à $20\,000$~M\$ d'actifs et
$15\,000$~M\euro{} de provisions. La méthode est complète, l'objet est fictif : le mémoire
démontre une faisabilité, il ne mesure pas une entité. C'est la première objection qu'un
rapporteur formule. Elle est juste. On la lève avec la seule donnée disponible sans accès
interne, les rapports SFCR de l'exercice 2024, qui sont des publications
réglementaires.

Trois grandeurs suffisent par entité : provisions techniques, fonds propres éligibles,
SCR. Chaque champ porte sa provenance, \emph{publié} quand il est lu tel quel,
\emph{déduit} quand il résulte d'une identité comptable ou d'un ratio de couverture publié.
Aucun n'est estimé. La conversion en dollars, nécessaire parce que le panel de calibration
exprime les actifs en \$, n'est pas un paramètre sensible : l'élasticité valant $0{,}074$, une
erreur de $10\,\%$ sur le cours déplace $\lambda$ de $0{,}7\,\%$.

\begin{attention}
\textbf{Deux précautions de vocabulaire, avant tout chiffre.} D'abord, il n'existe \emph{aucun}
module DORA dans la Formule Standard : la grandeur calculée ici est un besoin de
capital ORSA, donc de pilier~2, conformément à ce que pose le
chapitre~\ref{chap:reglementaire}. Écrire \og le SCR DORA de telle entité\fg{} ferait
croire à un objet réglementaire qui n'existe pas. Ensuite, le rapport de ce besoin au SCR
publié est une mise à l'échelle et non une part : diviser un besoin de pilier~2 par le
SCR réglementaire ne fait pas du premier une composante du second. Le rapport répond à
\og est-ce gros ou petit pour cette entité\fg{}, et à rien d'autre. Enfin les entités sont désignées
par leur classe de taille, non par leur nom : le modèle n'utilise que le total de bilan et les
provisions, jamais l'identité, et il serait donc gratuit d'accoler un nom de société à un état de
conformité \emph{supposé}. Les identités et les références SFCR exactes sont tenues hors
du mémoire, dans le code de calcul, qui en imprime la table de correspondance à chaque exécution.
% SOURCES-SCRIPTS: 65
\end{attention}

\begin{center}\small
\renewcommand{\arraystretch}{1.25}
\begin{tabular}{@{}l r r r r r@{}}
\toprule
\textbf{Entité} & \textbf{Actifs} & \textbf{SCR publié} & \textbf{Besoin ORSA} &
\textbf{Bande} & \textbf{Rapport} \\
 & M\euro & M\euro & M\euro & M\euro & \\
\midrule
Assureur non-vie A & $2\,319$ & $425$ & $112$ & $[88\,;122]$ & $26{,}4\,\%$ \\
Assureur non-vie B & $5\,305$ & $812$ & $136$ & $[104\,;146]$ & $16{,}7\,\%$ \\
Assureur vie C & $37\,845$ & $1\,587$ & $196$ & $[149\,;213]$ & $12{,}3\,\%$ \\
\textbf{Assureur vie D} & $\mathbf{309\,800}$ & $\mathbf{14\,800}$ & $\mathbf{294}$ &
$\mathbf{[233\,;316]}$ & $\mathbf{2{,}0\,\%}$ \\
\midrule
\emph{Entité notionnelle} & $19\,231$ & --- & $169$ & $[132\,;183]$ & --- \\
\bottomrule
\end{tabular}
\end{center}

\vspace{2pt}
\noindent \textbf{Les quatre SCR du dénominateur sont publiés, et cela a été vérifié
pièce par pièce.} Les deux assureurs vie donnent leur SCR en synthèse ; les deux
assureurs non-vie le donnent dans leur état quantitatif \textsc{s.25.01}, à $425$ et
$812{,}5$~M\euro. Aucun n'est reconstitué à partir d'un taux de couverture, si bien que la
colonne des rapports repose entièrement sur des grandeurs lues. Ce qui reste \emph{déduit} est
le total d'actif de trois entités, obtenu par l'identité comptable provisions plus fonds
propres, et l'approximation y est conservatrice : les autres passifs manquants
minorent l'actif, donc $\lambda$, donc le capital. Sur la quatrième entité, où le total d'actif
est publié, l'identité se vérifie au million près.

\vspace{2pt}
\noindent \textbf{Un résultat en sort, sur le repère forfaitaire lui-même.} Les deux états
\textsc{s.25.01} publient le module de risque \emph{opérationnel} de Formule Standard, ce qui
permet de confronter le repère de $3\,\%$ des provisions employé dans ce chapitre à la charge
effectivement calculée : il vaut $54{,}1$ contre un module publié de $59$ sur la première entité
non-vie, et $62{,}1$ contre $39{,}7$ sur la seconde, soit des écarts de $-8\,\%$ et de
$+56\,\%$. L'écart ne vient pas du règlement mais de l'approximation : le repère ne retient que la
branche provisions de l'équation~\eqref{reg:eq:scrop}, sans la branche primes dont la formule prend
le maximum, sans le plafond assis sur le BSCR, et sur une assiette qui mêle parfois provisions et
autres passifs. Il ne se cite donc que comme un ordre de grandeur de comparaison, jamais comme la
charge réglementaire ni comme l'une de ses bornes, et le module publié lui est préféré chaque fois
qu'il existe.
% SOURCES-SCRIPTS: 65

\vspace{4pt}
La première lecture est celle qu'on attendait. Sur la plus grande entité du panel, un assureur
vie de $310$~Md\euro{} d'actifs, le besoin de capital associé à la non-conformité vaut
$294$~M\euro{} en point et $[233\,;316]$ en bande, soit $2{,}0\,\%$ de son SCR publié :
un ordre de grandeur défendable, obtenu sans aucune donnée interne. Et l'écart entre états de
conformité subsiste à toutes les tailles, de $+41\,\%$ à $+49\,\%$, face à une charge forfaitaire
qui n'en distingue aucune. La thèse ne dépendait donc pas de l'entité fictive.

\begin{attention}
\textbf{La seconde lecture n'était pas attendue, et elle contraint le mémoire.} La charge décroît
beaucoup moins vite que la taille : un facteur $134$ sur les actifs ne divise la sévérité que par
$1{,}53$ et la fréquence que par $1{,}44$. Rapportée au SCR publié, la charge passe donc
de $2{,}0\,\%$ pour une grande entité à $26{,}4\,\%$ pour une petite. En posant qu'un seul risque
de non-conformité réglementaire ne peut raisonnablement peser plus de $10\,\%$ du SCR
total d'une entité, tous risques confondus, la bascule se situe entre $37\,845$ et
$309\,800$~M\euro{} d'actifs. Ce n'est pas un seuil fin, quatre entités ne le déterminent pas,
et il dépend aussi du levier propre de l'entité, dont le rapport SCR sur actifs va ici
de $4{,}2$ à $18{,}3\,\%$. C'est un ordre de grandeur, pour la cause que le
\S\ref{sec:ancrage} nommait sans la chiffrer : l'élasticité corrige l'\emph{échelle} de la
sévérité, jamais sa \emph{forme}, qui reste celle de grandes institutions financières.
\end{attention}

\begin{cle}
\textbf{Conséquence directe sur le chiffre central, qu'il faut assumer.} L'entité notionnelle
pèse $19\,231$~M\euro{} d'actifs et se trouve donc \emph{dans} la zone que la ligne précédente
déclare hors domaine, ce qui suffit à en tirer la conséquence : le $169$~M\euro{} doit être lu
comme une borne supérieure d'ordre de grandeur, non comme une charge plausible. Un
second test, celui du levier, ne l'accable pas : en imputant à cette entité le levier médian des
quatre entités réelles ($10{,}0\,\%$), soit $1\,932$~M\euro{} de SCR, la charge en
représenterait $9\,\%$, donc juste \emph{sous} le critère de $10\,\%$. La requalification repose
ainsi sur le seul argument de taille. Ce que le mémoire établit n'en est pas
affaibli, parce que sa thèse porte sur l'\emph{écart} entre états de conformité face à un
forfait insensible, et qu'un écart est un rapport : il survit à une erreur de niveau commune aux
trois états. Mais le \emph{niveau}, lui, cesse ici d'être présenté comme une mesure. Le lever
demanderait une sévérité calibrée sur des incidents d'entités de cette taille, c'est-à-dire
exactement la donnée que le chapitre~\ref{chap:donnees} déclare absente.
\end{cle}

Deux réserves enfin, à ne pas reléguer en note. Aucune de ces entités ne publie son état de
conformité DORA : on mesure une \emph{exposition} sous un état supposé, jamais une
non-conformité constatée : les trois états sont donc donnés côte à côte sans qu'aucun
soit désigné. Et deux des quatre entités publient \og provisions techniques et autres passifs\fg{}
sans les séparer, ce qui surestime l'assiette du repère forfaitaire pour elles ; pour les deux
assureurs vie, la branche provisions de la formule applique de surcroît un coefficient plus faible
que celui du repère, qui ne vaut donc pas pour elles. Le rapport au SCR publié, lui,
n'est pas concerné.

% SOURCES-SCRIPTS: 65

\section{Ce que l'agrégation réglementaire fait de cet écart}
\label{sec:agregation-scr}

Les sections qui précèdent chiffrent un besoin et un écart. Il reste à dire ce que le capital
total de l'entité en retient, et la réponse ne va pas de soi. Sous Solvabilité~II, le risque
opérationnel n'entre pas dans la matrice de corrélation du BSCR : il s'ajoute,
selon l'identité $\mathrm{SCR} = \mathrm{BSCR} + \mathrm{Adj} + \mathrm{SCR}_{Op}$ posée au
chapitre~\ref{chap:reglementaire}. Un besoin de nature opérationnelle porté en ORSA se
transmet donc au capital total sans aucun bénéfice de diversification, et cela vaut du
niveau comme de l'écart entre états. Un euro d'écart DORA coûte un euro de capital.

\paragraph{Le contrefactuel.} Pour mesurer cette convention, le même montant est agrégé comme un
module de plus, corrélé à $\rho$ avec le BSCR : un montant petit devant le BSCR
n'y est reconnu qu'à hauteur de $\rho$, quand la convention additive le reconnaît en entier. Sur
les quatre bilans réels, le BSCR n'étant que minoré, les taux lus sont des majorants
\matcomp.

\begin{center}\small
\renewcommand{\arraystretch}{1.25}
\begin{tabular}{@{}l r r r r@{}}
\toprule
\textbf{Entité} & \textbf{Écart C$\to$NC} & \textbf{$d = D/\mathrm{BSCR}$} &
\textbf{$\tau$ à $\rho = 0{,}25$} & \textbf{additif / corrélé} \\
 & M\euro & & & \\
\midrule
Assureur non-vie A & $34{,}9$ & $0{,}0954$ & $29{,}4\,\%$ & $3{,}41$ \\
Assureur non-vie B & $44{,}4$ & $0{,}0575$ & $27{,}7\,\%$ & $3{,}62$ \\
Assureur vie C & $62{,}7$ & $0{,}0514$ & $27{,}4\,\%$ & $3{,}65$ \\
Assureur vie D & $86{,}0$ & $0{,}0076$ & $25{,}4\,\%$ & $3{,}94$ \\
\bottomrule
\end{tabular}
\end{center}

\vspace{2pt}
\noindent À $\rho = 0{,}25$, la convention additive reconnaît ainsi de $3{,}4$ à $3{,}9$ fois ce
qu'un traitement corrélé retiendrait, et le rapport reste dans cette bande étroite sur un facteur
$134$ de taille de bilan, précisément parce que $\tau$ tend vers $\rho$. Lu sur le niveau plutôt
que sur l'écart, où $d$ est plus grand, le facteur va de $2{,}6$ à $3{,}8$. Rapporté au
SCR publié de chaque entité, l'écart entre états pèse de $0{,}58\,\%$ pour la plus grande
à $8{,}21\,\%$ pour la plus petite, et ces points de SCR sont acquis en entier.

\paragraph{Le plafond borne en outre ce qu'un forfait saurait exprimer.} Sur l'assureur non-vie~A,
le besoin calculé vaut $2{,}21$ fois la marge restante sous le plafond de $0{,}3 \times
\mathrm{BSCR}$ : même un forfait rendu sensible à la conformité saturerait avant de l'exprimer
\matcomp.

\begin{cle}
\textbf{La lecture décisionnelle, et elle renverse une intuition de gestion.} Un euro économisé
sur la non-conformité DORA vaut, en capital, davantage qu'un euro économisé sur un risque
qui se diversifie : le premier est reconnu en entier, le second à hauteur de sa corrélation. Une
direction des risques qui classerait ses chantiers de remédiation au seul vu de la sinistralité
attendue sous-estimerait donc le rendement en capital de celui-ci, d'un facteur de l'ordre de
trois à quatre. Cette conclusion ne dépend d'aucun paramètre du modèle de cascade : elle tient à
la place que le règlement assigne au risque opérationnel, et elle survivrait à un changement
complet de calibration.
\end{cle}

% SOURCES-SCRIPTS: 95 65

\section{Détenir ou transférer : le SCR confronté au prix de marché du risque}
\label{sec:detenir-transferer}

Les trois confrontations qui précèdent portent sur le \emph{niveau}. Il en manque une sur le
\emph{prix}. Un SCR n'est pas une perte : c'est du capital immobilisé, et le capital coûte.
Solvabilité~II chiffre ce coût explicitement par le taux de coût du capital de la marge de
risque, historiquement $6\,\%$ et ramené à $4{,}75\,\%$ par la directive
(UE)~2025/2, dont la transposition est due pour janvier 2027. Détenir le risque coûte donc
$\mathrm{CoC}\times\mathrm{SCR}$ par an. En face, le \emph{transférer} coûte une prime, que le
marché cyber facture avec un chargement : à un ratio de sinistralité de l'ordre de
$48{,}5\,\%$, l'assureur encaisse environ $2{,}1$ fois la prime pure. Les deux grandeurs sont
homogènes, donc comparables. Un directeur des risques fait cette comparaison.

\paragraph{Le multiple de capital, et ce qu'il coûte.} À l'échelle d'entité
(table~\ref{tab:echelle}, lecture composée), le modèle donne un SCR de $169$~M\euro{} pour une
perte attendue de $3{,}6$~M\euro, soit un multiple de capital de $47$ contre $8{,}3$
à l'échelle du secteur. L'écart est la mutualisation, lue à l'envers : le capital d'une entité
% SOURCES-SCRIPTS: 60
isolée est dominé par un sinistre unique, celui d'un secteur par l'accumulation. Le coût annuel
de détention vaut alors $8{,}0$~M\euro{} sous le régime révisé et $10{,}1$~M\euro{} sous
l'actuel, c'est-à-dire $2{,}2$ à $2{,}8$ fois la sinistralité attendue. Immobiliser
coûte plusieurs fois ce que le risque coûte en moyenne : c'est ce rapport qui rend le transfert
économiquement pertinent, et il ne se lit sur aucune des comparaisons de niveau.

\paragraph{Dimensionner plutôt que tout céder.} Un traité en excès de perte annuel de portée $L$
minimise le coût total, prime plus coût du capital résiduel, en $L^\star=\mathrm{SCR}=169$~M\euro{} :
le coût total passe de $8{,}0$ à $2{,}9$~M\euro{} par an, soit $-64\,\%$. Ce gain est une borne
supérieure, entamée par la capacité du marché, les exclusions et le risque de contrepartie, mais il
faudrait que $78\,\%$ de la couverture soit inefficace pour que la détention redevienne
préférable. Une entité devenue conforme voit en outre ce coût tomber de $81\,\%$, la conformité
déplaçant le point d'équilibre entre rétention et transfert. Le calcul, ses réserves et la figure
sont déportés \matcomp.
% SOURCES-SCRIPTS: 58 67

\begin{cle}
\textbf{La boucle que cette section referme.} DORA impose des exigences de \emph{gestion},
Solvabilité~II impose du \emph{capital}, et rien n'articule les deux. En passant par le prix,
le SCR cesse d'être un
chiffre à publier pour devenir un critère de décision entre trois leviers concurrents.
C'est la lecture que la thèse du mémoire appelait : le niveau est illustratif, mais le
\emph{rapport} entre le coût de détention et le prix de transfert, lui, est une grandeur
d'entité, robuste au changement d'échelle puisque ses deux termes se rapportent au même SCR.
\end{cle}

\begin{cle}
\textbf{Pourquoi le niveau ne peut pas être une mesure.} La confrontation ne calibre pas le
niveau, elle le
disqualifie comme mesure et confirme, par un chemin externe, la lecture que tout le
mémoire adopte : ce sont le rapport au socle, la largeur de la bande
d'ignorance et la hiérarchie des piliers qui portent les conclusions, jamais le
montant en euros.
\end{cle}

% =====================================================================
